% !PRINTED 262
\weberpage{262}
Es möge nun $A$ irgend ein Element der Gruppe $L_p$ sein,
und $U = \begin{pmatrix} \lambda, & \mu \\ \nu, & \varrho \end{pmatrix}$ ein Element aus $E_{2,p}$. Wir wollen aus $A$ durch
Transformation mit $U$ eine gewisse Normalform ableiten.

Es soll zunächst versucht werden, $A$ in die Normalform
\weq{1}{S = \begin{pmatrix} \sigma, & 0 \\ 0, & \sigma^{-1} \end{pmatrix}}
zu transformiren. Aus der Gleichung
\[
\begin{pmatrix} a, & b \\ c, & d \end{pmatrix} \begin{pmatrix} \lambda, & \mu \\ \nu, & \varrho \end{pmatrix} = \begin{pmatrix} \lambda, & \mu \\ \nu, & \varrho \end{pmatrix} \begin{pmatrix} \sigma, & 0 \\ 0, & \sigma^{-1} \end{pmatrix}
\]
erhält man
\weq{2}{\begin{aligned}
&(a-\sigma)\,\lambda + b\,\nu = 0, \quad (a-\sigma^{-1})\,\mu + b\,\varrho = 0,\\
&c\,\lambda + (d-\sigma)\,\nu = 0, \quad\ \ c\,\mu + (d-\sigma^{-1})\,\varrho = 0,
\end{aligned}}
und daraus ergiebt sich, dass $\sigma$ und $\sigma^{-1}$ die Wurzeln der quadratischen
Gleichung
\[
(a-\sigma)\,(d-\sigma) - b\,c = 0
\]
oder
\weq{3}{\sigma^2 - \sigma\,(a+d) + 1 = 0}
sein müssen. Hat man $\sigma$ und $\sigma^{-1}$ hieraus bestimmt, so findet
man aus (2) die Verhältnisse $\lambda : \nu$ und $\mu : \varrho$. Die Grössen $\lambda, \mu, \nu, \varrho$
selbst sind dann noch so zu bestimmen, dass $\lambda\varrho - \mu\nu = 1$ wird.
Dies ist immer möglich, wenn die aus (2) bestimmten Werthe
nicht die Gleichung $\lambda\varrho = \mu\nu$ erfüllen. Sind $b$ und $c$ beide $= 0$,
so hat $A$ schon die Form (1). Ist aber eine der beiden Zahlen
$b, c$, etwa $b$, von Null verschieden, so ergiebt sich aus (2):
\weq{4}{\frac{\nu}{\lambda} = -\frac{(a-\sigma)}{b}, \quad \frac{\varrho}{\mu} = -\frac{(a-\sigma^{-1})}{b},}
und es kann nur dann $\lambda\varrho = \mu\nu$ oder $\nu : \lambda = \varrho : \mu$ sein, wenn
$\sigma = \sigma^{-1}$, also $\sigma = \pm 1$ ist, d.\,h. nach (3), wenn $a + d = \pm 2$ ist.

Die Gleichung (3) ist aber im Körper $\mathfrak{E}_{2,p}$ immer lösbar, weil
jede \gesperrt{reelle} Zahl in diesem Körper ein Quadrat ist, und ergiebt
\weq{5}{\begin{aligned}
\sigma &= \frac{a+d}{2} + \sqrt{\left(\frac{a+d}{2}\right)^2 - 1}\\
\sigma^{-1} &= \frac{a+d}{2} - \sqrt{\left(\frac{a+d}{2}\right)^2 - 1}
\end{aligned}}
